\documentclass[lettersize,journal]{IEEEtran}
\usepackage[T1]{fontenc}
\usepackage{amsmath,amssymb,amsfonts,graphicx,booktabs,cite,url,textcomp}
\usepackage{tikz}
\usepackage[hidelinks]{hyperref}
\usetikzlibrary{arrows.meta,positioning,calc}
\graphicspath{{../python/latest/biotin_20260913/}}
\renewcommand{\thesection}{S\arabic{section}}
\renewcommand{\thefigure}{S\arabic{figure}}
\renewcommand{\thetable}{S\arabic{table}}
\renewcommand{\theequation}{S\arabic{equation}}
\newcommand{\mainpaper}[1]{\href{SPR_smart_stop_current.pdf}{main paper, #1}}
\begin{document}
\title{Supplementary Material: SPR Virtual Sensor for Early Stopping}
\author{Rodrigo T. Ara\'ujo, Mateus S. Marques,
Arthur A. Melo, and Antonio M. N. Lima
\thanks{\raggedright VIRTUS-CC, Department of Electrical Engineering, Universidade Federal de
Campina Grande, Campina Grande, Brazil. E-mail:
\{rodrigo.araujo,mateus.marques\}@ee.ufcg.edu.br;
\{arthur.aprigio,amnlima\}@dee.ufcg.edu.br.}
\thanks{\raggedright This work has been partially funded by SmartSens Base Systems,
VIRTUS-CC, with resources from PPI HardwareBR of MCTI, grant 055/2023,
signed with EMBRAPII.}}


\markboth{Supplementary Material: SPR Virtual Sensor}{Ara\'ujo \textit{et al.}}
\maketitle

\begin{abstract}
This supplement provides derivations and diagnostics beyond those presented
in the main article on the surface plasmon resonance virtual sensor.
The optical operator, moving-equilibrium error equation and parameter
ambiguity are developed from the main model without reproducing its
governing equations. Additional processing details specify centroid
integration and the fitted solution-exchange fronts. Reference-experiment
analyses provide transition-specific serial delays, fitted-prefix plots,
the complementary wavelength derivatives and normalization under alternative
time coordinates. For the independent experiment, additional wavelength
comparisons extend the measured-response panels in
the article. Endpoint sensitivity is documented for neighboring rinse
windows, vendor-angle processing and an alternative baseline; the primary
endpoint table remains in the main paper. Computational details specify
the residual-call budget, allocation measurement and reproduction files,
while the benchmark summary is cited from the article. Each section
identifies the corresponding main-paper material so that shared definitions,
results and experimental qualifications can be consulted directly. All
additional panels use the saved experimental observations and fits.
\end{abstract}
\begin{IEEEkeywords}
Delivery modeling, endpoint sensitivity, kinetic identifiability,
reproducibility, surface plasmon resonance.
\end{IEEEkeywords}
\raggedbottom
\noindent Read alongside the \mainpaper{Sections~II--V}. Notation,
experimental conditions and decision criteria follow that article.
Links below refer to its sections, equations, tables and figures.
\tableofcontents

\section{Additional model derivations}
\subsection{Optical operator and a general response target}
The forward chain and reference kinetic model are defined in the
\mainpaper{Eqs.~(1)--(3) and Fig.~1}; the conference extension is described
in its Introduction. An explicit form of the reflectance-minimum operator is
\begin{equation}
 H(d,n)=\arg\min_\theta\mathcal R_p(\theta;d,n)
       -\arg\min_\theta\mathcal R_p(\theta;0,n).
\end{equation}
This operator concerns the surface-layer contribution at a fixed bulk
condition; the bulk correction remains part of the main observation chain.
For a general remaining fraction $\varepsilon$, the constant-input solution
implies
\begin{equation}\label{eq:virtualtime}
 s_\varepsilon=-\tau\ln\varepsilon.
\end{equation}
The particular 99\% target and its injection-clock conversion are specified
in the \mainpaper{Eqs.~(6) and (12)}.

\subsection{Moving equilibrium under variable delivery}
Starting from the \mainpaper{Eq.~(4)}, let
$\lambda_j=1/\tau_{\rm rel,j}$,
$\Gamma_{\rm eq,j}=k_aC_j/\lambda_j$, and
$e_j=\Gamma_j-\Gamma_{\rm eq,j}$. Differentiating gives
\begin{equation}
 \dot e_j=-\lambda_j(t)e_j-\dot\Gamma_{\rm eq,j}(t).
\end{equation}
The additional forcing term shows why a fixed multiple of one time constant
cannot generally predict completion during solution exchange. A possible
delivery representation is
\begin{equation}
 C_j(t)=\int_0^\infty h_j(u)C_{\rm in}(t-u)\,du,
\end{equation}
where $h_j$ describes delay and dispersion. Applying this extension requires
a calibrated input and propagation under a declared future flow protocol;
the measured TIR fronts alone do not supply active-reagent concentration
(\mainpaper{Section~II-B}).

\subsection{Parameter ambiguity and forecast sensitivity}
The non-identifiability discussed in the \mainpaper{Section~II-D} can be
written explicitly. For any $0<\widetilde k_d<k_{\rm obs}$, the family
\begin{equation}
 \widetilde k_a=(k_{\rm obs}-\widetilde k_d)/C,\qquad
 \widetilde R_{\max}=R_{\rm eq}\frac{k_{\rm obs}}{k_{\rm obs}-\widetilde k_d}
\end{equation}
preserves the same constant-concentration association curve.
If $\widehat k_{\rm obs}=k_{\rm obs}(1+\delta)$ with $\delta>-1$,
substitution into \eqref{eq:virtualtime} gives
\begin{equation}
 \frac{\widehat s_{99}}{s_{99}}=\frac{1}{1+\delta},\qquad
 \frac{R_{\rm eq}-R(\widehat s_{99})}{R_{\rm eq}-R_0}
 =100^{-1/(1+\delta)}.
\end{equation}
This is an analytical sensitivity result, not a synthetic validation
experiment. The derivative-based remaining-response relation and its
limitations are given in the \mainpaper{Eq.~(7) and Section~II-F}.

\section{Additional extraction and front-fitting details}
\subsection{Centroid integration}
For the angular preprocessing and support-selection rule in the
\mainpaper{Section~III-D}, the centroid on the selected lobe $\mathcal L$ is
\begin{equation}
 x_c=\frac{\int_{\mathcal L}x[I_\ell-I(x)]\,dx}
 {\int_{\mathcal L}[I_\ell-I(x)]\,dx}.
\end{equation}
Boundary crossings are interpolated and integration uses the trapezoidal
rule. The tested TIR alternatives comprise derivative-peak,
derivative-centroid and template-shift estimators. Absolute TIR angles
extrapolated from the SPR calibration are not treated as an independently
calibrated observable. Bulk coefficients, baseline convention and the
noise statistic are specified in the \mainpaper{Section~III-D}; derivative
uncertainty and bootstrap settings are specified in its Sections~III-F--G.

\subsection{Local transition fit}
The timing procedure in the \mainpaper{Section~III-E} fits
\begin{equation}
 z(t)=b+m(t-t_c)+sA[1+e^{-(t-t_{50})/a}]^{-1},
\end{equation}
where $s$ specifies transition direction and $a>0$. The fitted markers are
$t_{10}=t_{50}-a\ln9$ and $t_{90}=t_{50}+a\ln9$.
For a decline, these denote fractions of downward exchange.
The linear background absorbs local drift. A separate affine-scaled,
time-shifted trace comparison provides a check independent of the logistic
shape. Section~S3 shows the individual transitions; the complete
sensorgrams remain in the \mainpaper{Fig.~2(A,B)}.

\section{Complementary reference-experiment diagnostics}
\subsection{Individual solution-exchange fronts}
Figure~\ref{fig:fronts} expands the transition regions from the
\mainpaper{Fig.~2(A)}. Table~\ref{tab:delay} gives the transition-specific
delays underlying the range reported in its Section~IV-B.
\begin{figure*}[tp]\centering
\includegraphics[width=.96\textwidth]{paper_figures/serial_fronts.pdf}
\caption{Individual normalized TIR transitions. Colored dotted lines mark
the fitted midpoint of each readout; black dashed lines mark the available
injection events. Normalization removes each fitted local background and
amplitude. These panels expand the front shapes without repeating the
full-record traces in the main paper.}
\label{fig:fronts}
\end{figure*}
\begin{table*}[t]\centering
\caption{Channel 1 minus channel 2 fitted $t_{50}$ (min).}
\label{tab:delay}
\begin{tabular}{lrr}\toprule
Transition & L1--L3, 670 nm & L2--L5, 785 nm\\\midrule
PBS rise & 2.215 & 2.221\\
PBS decline & 2.185 & 2.195\\
Biotin rise & 2.249 & 2.245\\
Water-associated decline & 2.236 & 2.256\\\bottomrule
\end{tabular}
\end{table*}

\subsection{Fitted prefixes and complementary derivative readouts}
The parameter estimates are reported in the \mainpaper{Table~2}.
Figure~\ref{fig:prefixfits} supplies the measured-prefix fit and later
continuation for each readout. Figure~\ref{fig:extra-slopes} adds the
785-nm slope panels; the corresponding 670-nm panels and summary statistics
are in the \mainpaper{Fig.~2(C,D) and Table~1}.
\begin{figure*}[tp]\centering
\includegraphics[width=.96\textwidth]{paper_figures/tail_fits.pdf}
\caption{Reference-run fitted prefixes and subsequent continuation for all
four corrected readouts. The vertical line separates fitting and prediction
intervals; the horizontal line is the fitted asymptote. The water-associated
decline is excluded from the association target. Fit values and intervals
are given in the \mainpaper{Table~2}.}
\label{fig:prefixfits}
\end{figure*}
\begin{figure*}[tp]\centering
\includegraphics[width=.96\textwidth]{paper_figures/supplement_slopes_785.pdf}
\caption{Additional 785-nm derivative readouts, L2 and L5. Shading denotes
the approximate corrected-slope 95\% interval; dotted lines mark
$\pm0.5$ mdeg/min. These are the wavelength panels absent from the
\mainpaper{Fig.~2(C,D)}.}
\label{fig:extra-slopes}
\end{figure*}

\subsection{Timing decomposition and window sensitivity}
Using the fitted values from the \mainpaper{Table~2}, define
\begin{equation}
 u_j=\frac{T-T_{0,j}}{\widehat\tau_j},\qquad
 z_j=\frac{R_j-b_j}{A_j}.
\end{equation}
Table~\ref{tab:timing} separates the fitted tail-start offset from relaxation;
Fig.~\ref{fig:normalization} displays the alternative clocks. Interpretation
of the shared dimensionless target and the equal-local-age refit is given
in the \mainpaper{Section~IV-E}. For that refit, the respective RMS residuals
are 1.00, 0.68, 0.92 and 0.66 mdeg for L1, L2, L3 and L5. Small residuals
therefore coexist with the rate sensitivity reported in the article.
\begin{table*}[tp]\centering
\caption{Decomposition of the main-paper forecasts, in minutes.
$T_0$ is the causal tail-fit start and $a_{50}$ the retrospective local TIR
midpoint; they are different timing markers.}
\label{tab:timing}
\begin{tabular}{lrrr}\toprule
Readout & $T_0$ & $4.605\tau$ & $T_{99}-a_{50}$\\\midrule
L1 & 7.17 & 9.30 & 9.75\\
L2 & 7.17 & 8.82 & 9.28\\
L3 & 4.46 & 19.47 & 19.45\\
L5 & 4.46 & 18.25 & 18.24\\\bottomrule
\end{tabular}
\end{table*}
\begin{figure*}[tp]\centering
\includegraphics[width=.96\textwidth]{normalization_results/normalized_kinetics.pdf}
\caption{Reference response under injection, local-arrival and fitted-time-
constant normalization, with the forecast decomposition. Collapse of fitted
curves follows the exponential model; it is not independent evidence of
identical molecular dynamics.}
\label{fig:normalization}
\end{figure*}

\section{Additional repeat-experiment and endpoint diagnostics}
\subsection{Complementary measured-response comparisons}
The validation experiment tests the approximately 25-min duration selected
from the reference forecast, as specified in the
\mainpaper{Sections~III-A--C and IV-F--H}. Its role is to provide the
experimental outcome at that duration, not to generate a new prediction.
Figure~\ref{fig:extra-repeat} adds L2 and L5 to the measured L1/L3 comparison
in the main paper. No kinetic fit or stopping criterion is applied to
these validation traces.
\begin{figure*}[tp]\centering
\includegraphics[width=.96\textwidth]{paper_figures/supplement_short_785.pdf}
\caption{Additional 785-nm measured-response comparisons: L2 and L5 during
exposure and after rinse. Baseline, reference-prediction and command
markers, and rinse alignment follow the \mainpaper{Fig.~3}. No amplitude
rescaling or validation-run model fit is applied.}
\label{fig:extra-repeat}
\end{figure*}

\subsection{Additional endpoint windows and baseline conventions}
The exposure definition and primary retained-response comparison are in the
\mainpaper{Eq.~(13) and Table~3}. The tables below provide the neighboring
rinse windows, vendor-angle sensitivity and initial-water baseline.
The main-paper primary centroid rows are omitted. Values are within-run
temporal means, in mdeg except ratios; rinse age follows each readout's
local TIR midpoint. These comparisons use the qualification and limitations
stated in the \mainpaper{Section~IV-H}.
\begin{table*}[htbp]\centering\small
\caption{Additional rinse windows: centroid, pre-injection baseline. The primary 15--20 min window is reported in the \mainpaper{Table~3}.}
\begin{tabular}{llrrrr}\toprule
Readout & Rinse age (min) & Reference & Repeat & Difference & Ratio (\%)\\\midrule
L1 & 10--15 & 35.19 & 23.05 & -12.14 & 65.5\\
L1 & 20--24 & 29.98 & 21.22 & -8.76 & 70.8\\
L2 & 10--15 & 21.19 & 14.09 & -7.10 & 66.5\\
L2 & 20--24 & 18.33 & 12.94 & -5.39 & 70.6\\
L3 & 10--15 & 34.47 & 9.07 & -25.41 & 26.3\\
L3 & 20--24 & 27.88 & 6.42 & -21.45 & 23.0\\
L5 & 10--15 & 19.81 & 5.32 & -14.49 & 26.9\\
L5 & 20--24 & 16.76 & 3.86 & -12.90 & 23.1\\
\bottomrule\end{tabular}
\end{table*}

\begin{table*}[htbp]\centering\small
\caption{Vendor, pre-injection baseline.}
\begin{tabular}{llrrrr}\toprule
Readout & Rinse age (min) & Reference & Repeat & Difference & Ratio (\%)\\\midrule
L1 & 10--15 & 34.27 & 21.46 & -12.81 & 62.6\\
L1 & 15--20 & 30.10 & 18.77 & -11.32 & 62.4\\
L1 & 20--24 & 28.91 & 19.32 & -9.59 & 66.8\\
L2 & 10--15 & 20.38 & 12.50 & -7.88 & 61.3\\
L2 & 15--20 & 18.42 & 11.17 & -7.25 & 60.6\\
L2 & 20--24 & 17.60 & 11.33 & -6.26 & 64.4\\
L3 & 10--15 & 35.13 & 9.25 & -25.88 & 26.3\\
L3 & 15--20 & 31.46 & 7.48 & -23.98 & 23.8\\
L3 & 20--24 & 28.17 & 6.37 & -21.80 & 22.6\\
L5 & 10--15 & 20.34 & 5.25 & -15.08 & 25.8\\
L5 & 15--20 & 18.60 & 4.35 & -14.24 & 23.4\\
L5 & 20--24 & 17.46 & 3.75 & -13.71 & 21.5\\
\bottomrule\end{tabular}\end{table*}

\begin{table*}[htbp]\centering\small
\caption{Centroid, initial-water baseline.}
\begin{tabular}{llrrrr}\toprule
Readout & Rinse age (min) & Reference & Repeat & Difference & Ratio (\%)\\\midrule
L1 & 10--15 & 28.52 & 40.46 & +11.95 & 141.9\\
L1 & 15--20 & 24.52 & 38.09 & +13.56 & 155.3\\
L1 & 20--24 & 23.31 & 38.64 & +15.32 & 165.7\\
L2 & 10--15 & 17.64 & 23.55 & +5.91 & 133.5\\
L2 & 15--20 & 15.59 & 22.27 & +6.68 & 142.9\\
L2 & 20--24 & 14.78 & 22.40 & +7.62 & 151.6\\
L3 & 10--15 & 47.78 & 33.32 & -14.46 & 69.7\\
L3 & 15--20 & 44.16 & 31.37 & -12.78 & 71.0\\
L3 & 20--24 & 41.18 & 30.68 & -10.51 & 74.5\\
L5 & 10--15 & 27.25 & 19.66 & -7.59 & 72.1\\
L5 & 15--20 & 25.49 & 18.73 & -6.77 & 73.5\\
L5 & 20--24 & 24.21 & 18.20 & -6.00 & 75.2\\
\bottomrule\end{tabular}\end{table*}




\section{Computational accounting and reproduction details}
\subsection{Replay configuration and residual-call budget}
Benchmark hardware, summary timings, operation totals and allocation sizes
are reported in the \mainpaper{Section~V-A and Table~4}.
The common eligible interval spans 11.22--21.50 min after injection.
The environment used Python 3.12.14, NumPy 2.5.3 and SciPy 1.18.1.
For the residual arithmetic defined in that section, instrumenting all
callbacks, including those used by the numerical Jacobian, gives
\begin{equation}
 F_{\rm residual}=4\sum_{j=1}^{4}N_jQ_j,\qquad
 N_{\rm expm1}=\sum_{j=1}^{4}N_jQ_j,
\end{equation}
where $N_j$ is the current tail size and $Q_j$ is the residual-call count
across all four initializations. No fixed FLOP equivalent is assigned to
the transcendental operation. Negations are tracked separately in the
export; the solver and covariance work are accounted for by elapsed time,
not by this residual-only sum.
The current implementation uses four initial rates and
\texttt{max\_nfev}=1500 per start. This limit applies to objective
evaluations separately from numerical-Jacobian evaluations.
It is not a 40-iteration recursive solver. For three fitted parameters and
a two-point Jacobian, the residual-call budget is at most
$4(1+3)1500$ per readout; conditional on a buffer cap $N_{\max}$, this gives
\begin{equation}
 F_{\rm residual}\leq
 4\cdot4\cdot4\cdot(1+3)\cdot1500\,N_{\max}.
\end{equation}
This bounds only residual arithmetic under that configuration, not total
computation. The present prefixes grow with acquisition: 127 is the
observed maximum, not an implemented permanent memory cap. A deployment
requiring a fixed worst-case budget must impose and validate such a cap.


\subsection{Allocation measurement}
The payload reported in the \mainpaper{Table~4} counts float64 tail times
and responses. The incremental \texttt{tracemalloc} peak was measured
during the final update after preloading. The shortest recorded acquisition
interval was 8.022 s. Scope exclusions and the distinction between desktop
timing and embedded feasibility are specified in the
\mainpaper{Section~V-A}.
\subsection{Reproduction files and consistency checks}
The analysis repository contains the processing and numerical exports
under \path{python/latest/biotin_20260913} and
\path{python/latest/biotin_20260914}. The validation-experiment results use the processed-response CSV files
and measured front/endpoint comparisons; kinetic fits and stopping replays
of that experiment are excluded from this article. Reproduction of its
figures requires only saved measured responses and front timing, not
estimated kinetic parameters. All 672 scans per selected readout passed
the extraction quality checks. The reference experiment alone supplies
the reported kinetic forecasts and decision analysis.
The timing script \path{benchmark_streaming.py} exports all 231 measured
latencies and a separately instrumented residual-call replay. All 77
callback counts agree with objective evaluations plus three times the
numerical-Jacobian evaluation count. These checks concern implementation
consistency, not chemical validity or between-sensor equivalence.


\end{document}
